\subsection{Pointwise convergence on the original big cell}
\label{archimedean-big-cell-continuity}

The almost-everywhere boundary in the preceding norm argument can be
removed on the regular big cell. The key estimate concerns the
original unipotent variables, rather than an abstract representative
of a radial $L^2$ function.

\begin{lemma}\label{block-gauss-polynomial-properness}
Let $n=2r+k+1$, $k\ge1$, and $N$ be the original flag unipotent
subgroup. Let $w$ exchange each outer isotropic vector with its
paired opposite and fix the middle space. For
$c=(c_0,\ldots,c_{r-1})\in(\mathbb C^\times)^r$ and
$\mu\in U(k,1)$ write
\[
 g(c,\mu)=wD(c,\mu),\quad
 D(c,\mu)=
 \operatorname{diag}(c_0,\ldots,c_{r-1},\mu,
                    \bar c_{r-1}^{-1},\ldots,\bar c_0^{-1}).
\]
For $c$ in a fixed compact subset of $(\mathbb C^\times)^r$,
there is a constant $C$ such that, for all $n_L,n_R\in N$,
\[
 \max\{\|n_L\|_{\rm HS},\|n_R\|_{\rm HS}\}
 \le C\|n_L^{-1}g(c,\mu)n_R\|_{\rm HS}^{\,n-1}.
\]
The constant is independent of $\mu$. The original compact
stabilizers can be included on the two outside factors by changing
$C$ by a fixed finite factor.
\end{lemma}

\begin{proof}
Use the ordered block sizes $(1^r,k+1,1^r)$.
The defining flag conditions and Hermitian invariance make $N$
upper block-unitriangular with identity middle block.
Hermitian invariance gives the opposite flag conditions, so there
are no omitted entries below this block diagonal.
Conjugation by $w$ makes it lower block-unitriangular.
For $X=n_L^{-1}g(c,\mu)n_R$ we therefore have the exact
block Gauss decomposition
\[
             w^{-1}X=LDU,\qquad
 L=w^{-1}n_L^{-1}w,\quad U=n_R,\quad D=D(c,\mu).
\]

Let $J$ be an initial segment ending at a block boundary before
the last block, and let $\Delta_J$ be the determinant of the
corresponding principal submatrix of $w^{-1}X$.
It equals the determinant of $D_J$.
The permitted boundaries are
\[
                1,\ldots,r,\quad r+k+1,\ldots,n-1 .
\]
There is no boundary inside the middle $(k+1)$-dimensional block.
Since $|\det\mu|=1$, the absolute values $|\Delta_J|$
are products of $|c_a|$ and their reciprocals.
They have a common positive lower bound $\delta$ on the specified
compact set; decrease $\delta$ to at most one.

Put $A=w^{-1}X$. If $i$ belongs to the last block of $J$ and
$t$ lies after $J$, then
\[
 U_{it}=(A_J^{-1}A_{J,t})_i,\qquad
 L_{ti}=(A_{t,J}A_J^{-1})_i .
\]
Indeed, multiplying the block factors cancels $L_JD_J$ in the
first expression and $D_JU_J$ in the second; the last block row
of $U_J^{-1}$ and last block column of $L_J^{-1}$ are the
corresponding rows and columns of the identity.
Cramer's rule expresses each entry as a determinant of size
$|J|\le n-1$ divided by $\Delta_J$.
Hadamard's inequality bounds its numerator by
$\|A\|_{\rm HS}^{|J|}$.
Since $\|X\|_{\rm HS}\ge1$ and $w$ is Euclidean unitary, this gives
\[
 \|L\|_{\rm HS},\|U\|_{\rm HS}
       \le n\delta^{-1}\|X\|_{\rm HS}^{n-1}.
\]
The original Hermitian matrix is a unitary involution, so
$n_L^{-1}=Hn_L^*H$ and
$\|n_L^{-1}\|_{\rm HS}=\|n_L\|_{\rm HS}$.
This proves the bound. For $r=0$, $N=\{1\}$ and it holds directly.

Each original compact stabilizer acts only on the middle block.
It normalizes $N$, commutes with $w$ and the outer diagonal
entries, and its conjugation distorts the Hilbert--Schmidt norm
by a bounded factor. Move the two compact factors into these
conjugations and replace $\mu$ by their left-right translate.
The latter still lies in $U(k,1)$. This proves the last assertion
without changing any Haar measure.
\end{proof}

\begin{theorem}\label{original-big-cell-continuity}
Let $F$ be a continuous original operator-valued coefficient with
\[
                  \|F(g)\|_{\rm HS}\le C_F\|g\|_{\rm HS}^{-A},
             \qquad A>4n^2(n-1).
\]
Its original two-sided Bessel integral is absolutely convergent
at every big-cell point. For every original smooth section,
the integral is continuous on
\[
 \mathcal Q=(\mathbb C^\times)^r
             \times\{\zeta\in\mathbb C:|\zeta|^2>TT'\}.
\]
Both assertions hold for every original unitary character scale.
\end{theorem}

\begin{proof}
Keep the original compact, unipotent and character factors in
the integral. The compact representation factors and the
characters have operator norm one.
For $c$ in a compact set away from zero, the preceding lemma
majorizes the norm of its integrand by
\[
 C'\bigl(\max\{1,\|n_L\|_{\rm HS},\|n_R\|_{\rm HS}\}\bigr)^{-A/(n-1)} .
\]
The ambient packing argument of
Section~\ref{archimedean-quotient-mass} applies to the closed
unimodular subgroup $N$: its Haar volume in a matrix ball of radius
$R\ge1$ is at most $C_NR^{2n^2}$.
Consequently the product Haar volume in the displayed maximum ball
is at most $C_N'C_NR^{4n^2}$, where the two constants retain their
possibly different original Haar masses.
Summing over dyadic annuli proves integrability, because
$A/(n-1)>4n^2$. The compact measures are finite.
No probability normalization or representative change has been
made in the integral.

The original section is smooth on the displayed exterior domain
\cite{big-cell-continuity-source}. Thus the integrand is continuous
in its big-cell parameters for each fixed integration variable.
On every compact subset of $\mathcal Q$ the same integrable
majorant applies. Dominated convergence gives continuity of the
actual integral, with its prescribed values.
\end{proof}

For the family used above, this theorem has no additional gap cost.
Set $q=r+1$ and let $h$ be its original noncompact gap.
The common effective gap $h/4$ gives
$A=h/(32\sqrt{2q})-1$. The original local hypotheses give
\[
 A\ge \frac{C_g}{2}q^4(3qn^2+n^3)-1
       \ge 500000n^3-1>4n^2(n-1).
\]
The smooth coefficient estimate of
Theorem~\ref{original-smooth-decay-from-tensor} supplies this
decay for every coefficient in the finite shifted family.
The constants in the integrable majorant are used only for
convergence and continuity; none enters the uniform power of $b$.

\begin{corollary}\label{original-pointwise-pieri-expansions}
Under the exact hypotheses of
Theorem~\ref{original-iterated-minor-norm}, its finite coefficient
expansions hold at every point of the original big cell.
With the source's normalized coefficient representatives, this
supplies $\mathrm{BesselExpansionNearParameter}$ and
$\mathrm{BesselExpansionCoefficientSum}$, with the same constants,
support bound and original real parameter reserve.
\end{corollary}

\begin{proof}
Fix any admissible datum and the finitely many prescribed
coefficient representatives. Every term in its expansion,
and the original coefficient multiplied by the scalar leading
minors, is continuous on $\mathcal Q$ by the theorem.
The preceding subsection proves equality almost everywhere for
the original radial measure. That measure has positive density
throughout $\mathcal Q$:
\[
 \prod_{a<r}|c_a|^{2(2r+k-2a-1)}
                    (|\zeta|^2-TT')^{k-1}\,dc\,d\zeta ,
\]
up to its retained positive normalization.
A nonzero continuous difference would be bounded away from zero
on a nonempty open set of positive measure. Thus the difference
vanishes at every point.

The argument applies to the arbitrary normalized representatives
quantified by the source. A representation whose entire Jacquet
image is zero contributes no coefficient; continuity also makes
that vanishing pointwise. The absolute coefficient sum remains
at most $b^{C_bm^2|\alpha|}$ and its support remains
$\|\eta\|_1\le(r+1)|\alpha|$, as already proved.
No uniform exceptional set over changing arithmetic data is needed:
the conclusion is pointwise for each admissible datum.
\end{proof}

This removes the big-cell representative-identification issue.
It does not yet supply a uniform point-evaluation bound relative
to the Jacquet norm. It also does not compare the original
orbital measures on lower-dimensional relevant strata, where a
block boundary minor used in the Gauss argument can vanish.
Those strata occur in the rational-return problem and remain
part of the weighted orbital frontier.
